%! Author = denis
%! Date = 11/08/2026

\begin{frame}{Mapa de Karnaugh}
    \small

    \begin{block}{Definição}
        O \textbf{Mapa de Karnaugh} é um método gráfico utilizado para
        simplificar expressões booleanas e, consequentemente,
        reduzir a quantidade de portas lógicas de um circuito.
    \end{block}

    \begin{columns}[T]

        \begin{column}{0.49\textwidth}
            \begin{block}{Objetivo}
                O método permite:

                \begin{itemize}
                    \item Representar uma tabela verdade;
                    \item Agrupar combinações semelhantes;
                    \item Eliminar variáveis;
                    \item Obter uma expressão simplificada.
                \end{itemize}
            \end{block}
        \end{column}

        \begin{column}{0.49\textwidth}
            \begin{block}{Ideia principal}
                Os valores $1$ são agrupados em grupos de:

                \[
                    \boxed{1,\;2,\;4,\;8,\;16,\ldots}
                \]

                Quanto maior o grupo, maior poderá ser a
                simplificação.
            \end{block}
        \end{column}

    \end{columns}

\end{frame}